% Generated draft: requires mathematical and source review.
\documentclass[11pt]{article}
\usepackage[margin=25mm]{geometry}
\usepackage{fontspec}
\setmainfont{Latin Modern Roman}
\usepackage{amsmath,amssymb,mathtools,unicode-math}
\setmathfont{Latin Modern Math}
\usepackage{enumitem,needspace,xurl,hyperref}
\setlength{\parindent}{0pt}
\setlength{\parskip}{5pt}
\setlength{\emergencystretch}{3em}
\newcommand{\N}{\mathbb{N}}
\newcommand{\Z}{\mathbb{Z}}
\newcommand{\Q}{\mathbb{Q}}
\newcommand{\R}{\mathbb{R}}
\newcommand{\C}{\mathbb{C}}
\newcommand{\F}{\mathbb{F}}
\newcommand{\A}{\mathbb{A}}
\newcommand{\PP}{\mathbb{P}}
\newcommand{\E}{\mathbb{E}}
\newcommand{\eps}{\varepsilon}
\newcommand{\dd}{\,\mathrm{d}}
\begin{document}
\section*{UnsolvedMath ID: 30003043}
\textbf{Source problem code:} catalogue entry
\subsection{Definitions and mathematical statement}
Let \(m\geq 18\) be an even integer and put \(r=\lfloor m/4\rfloor\). For a quadratic form \(\varphi\) on an \(m\)-dimensional vector space \(V\) over a field \(F\), write \(X_r(\varphi)\) for the orthogonal Grassmannian parameterizing \(r\)-dimensional totally isotropic subspaces: for a field extension \(E/F\), its \(E\)-points correspond to subspaces \(U\subseteq V\otimes_F E\) with \(\dim_E U=r\) and \(\varphi_E|_U=0\). The usual smooth projective interpretation requires appropriate nonsingularity assumptions, which the supplied record does not explicitly specify.

For a nonempty proper variety \(X/F\), define its index by
\[
\operatorname{ind}(X)=\gcd\{[\kappa(x):F]:x\text{ is a closed point of }X\},
\]
where \(\kappa(x)\) is the residue field of \(x\).

The source uses \(I_q^3\) for a condition in the fundamental-ideal filtration of quadratic Witt theory. In characteristic different from \(2\), the usual interpretation is that the Witt class \([\varphi]\) belongs to \(I(F)^3\), where \(W(F)\) is the Witt ring and \(I(F)\) is its fundamental ideal of even-dimensional classes. The supplied record gives neither a characteristic restriction on \(F\) nor a characteristic-dependent definition of \(I_q^3\). It also gives no definition of a generic form or of the family relative to which genericity is taken. These qualifications remain unresolved and must be fixed from the source before the assertion has a fully specified scope.

\textbf{Conjecture as supplied.} For every even integer \(m\geq18\), every field \(F\), and every generic \(m\)-dimensional quadratic form \(\varphi/F\) in \(I_q^3\), with the admissible fields, genericity, and quadratic-form conventions understood as intended by the source,
\[
\operatorname{ind}\bigl(X_{\lfloor m/4\rfloor}(\varphi)\bigr)
=2^{\lfloor m/4\rfloor}.
\]
The quantification over fields is subject to the unresolved source conventions just described; this display does not assert that every characteristic is admissible.
\subsection{Short English statement}
The supplied conjecture says that, for every even \(m\geq18\), the variety of \(\lfloor m/4\rfloor\)-dimensional totally isotropic subspaces of a generic \(m\)-dimensional quadratic form in \(I_q^3\) has index exactly \(2^{\lfloor m/4\rfloor}\). The field restrictions, nonsingularity assumptions, and meaning of genericity require clarification from the source.
\subsection{Sources}
\begin{itemize}
\item Supplied local UnsolvedMath registry record, permanent ID \(30003043\), title \emph{Exact Indices of Orthogonal Grassmannians}, exact source problem code \texttt{catalogue entry}. The record gives proposed year \(2016\) but names no individual proposer.
\item Supplied primary-source citation: \emph{Algebraic Cobordism and Projective Homogeneous Varieties} (2016), Oberwolfach Reports, pp.\ 215--265, DOI \texttt{10.4171/owr/2016/5}. Source URL and literature-source link labeled ``original OWR report'': \url{https://doi.org/10.4171/owr/2016/5}. No precise page or named contributor for this conjecture is supplied.
\item Supplied literature-source link labeled ``Karpenko on 16-dimensional forms'': \url{https://sites.ualberta.ca/~karpenko/publ/a16i3.pdf}. The record supplies no further bibliographic details for this link.
\item Literature assessment dated exactly \texttt{2026-08-21}: the supplied curation labels the problem open, describes the exact index formula as explicitly conjectural in the primary source, and states that related work on \(I\)- and quadratic-form invariants does not verify it. This is supplied dated curation, not an independently verified literature assessment. The background expressly requests independent verification of time-sensitive status and cited claims.
\end{itemize}
\end{document}
